\partstart{Tổng kết và hướng phát triển}
\subsection{Các mục tiêu đã thực hiện}
Đề tài đã xây dựng quy trình phân loại nguy cơ giao trễ trên DataCo,
có kiểm toán schema, nhãn, missing, outlier và cấu trúc đơn nhiều dòng.
Allowlist loại các trường kết quả và định danh; chia train/test và CV
theo mã đơn hạn chế rò rỉ giữa các dòng chung đơn.

K-NN và cây quyết định được so sánh trên A/B với cây mặc định và Dummy.
Mô hình được chọn bằng F1 CV; tiền xử lý fit trong từng fold, test chỉ
đánh giá sau khi khóa lựa chọn. Cấu hình, hash, phiên bản, bảng và hình
được lưu để truy vết.

Ma trận nhầm lẫn, khoảng cách train--CV, ablation lịch, metric theo đơn
và lỗi theo dịch vụ bổ sung cho đánh giá tổng thể, chỉ ra điểm yếu
Standard Class thay vì chỉ xếp hạng thuật toán.

\subsection{Kết quả chính và mức độ kết luận}
Trên mẫu 5.999 dòng, cây B được chọn với độ sâu 4, tối thiểu 20 mẫu/lá,
gini và ba đặc trưng lịch bổ sung. CV F1 đạt 0,7058, test F1 theo dòng
0,7098; Precision 0,8533 và Recall 0,6076. Các giá trị này mô tả kết quả
của đúng một protocol và một tập con, không phải đánh giá trên toàn bộ
180.519 dòng hay một dữ liệu ngoài độc lập.

K-NN B cải thiện so với A trong kết quả đã lưu, nhưng F1 train gần 1 và
CV thấp hơn đáng kể cho thấy cần thận trọng về tổng quát hóa. Cây baseline
cũng có khoảng cách train--CV lớn. Cây B nông có gap nhỏ hơn, song bỏ sót
248/254 dòng dương Standard Class. Vì vậy, cải thiện về kiểm soát độ phức
tạp chưa giải quyết được thất bại ở nhóm dịch vụ này.

Dummy đạt F1 dòng 0,6991 và F1 đơn 0,7123 bằng cách dự đoán mọi trường
hợp dương. Cây B đạt F1 đơn 0,7100, với Precision cao hơn nhưng Recall
thấp hơn nhiều. Kết quả nhấn mạnh cần đánh giá baseline và đơn vị nghiệp
vụ ngay từ thiết kế, thay vì coi một ngưỡng F1 tuyệt đối là đủ để kết
luận mô hình tốt. Hiện chưa có bằng chứng đủ để khuyến nghị tự động ra
quyết định vận chuyển.

\subsection{Điểm mạnh và hạn chế của quy trình}
Điểm mạnh là kiểm soát rò rỉ ở nhiều tầng: đặc trưng, nhóm và nơi fit
tiền xử lý. Sự nhất quán giữa cấu hình, runner và notebook giúp giảm sai
khác protocol. Manifest và artifact tổng hợp tạo đường truy vết từ số
liệu đến bảng/hình báo cáo. Các ví dụ lỗi được lọc thông tin cá nhân,
trong khi dữ liệu thô và model cục bộ không cần đưa vào bản công bố.

Hạn chế quan trọng của dữ liệu là nhãn âm chứa đơn hủy, thời điểm nghiệp
vụ của một số cột chưa có xác nhận độc lập và nguồn có thể không đại diện
cho vận hành tương lai. Hạn chế thực nghiệm là lấy một mẫu nhỏ, một seed,
ba fold và một holdout nhóm ngẫu nhiên. Chưa có khoảng tin cậy theo đơn,
đánh giá ngoài thời gian hoặc trên doanh nghiệp khác.

Hạn chế mô hình gồm khoảng cách phụ thuộc mã hóa ở K-NN, cây nhạy với
mẫu và nhóm hiếm, xác suất chưa hiệu chuẩn và các ngưỡng chưa gắn chi
phí. Không gian tìm kiếm nhỏ cũng giới hạn phạm vi kết luận về siêu
tham số. Những điểm chưa thực hiện không được mô tả như tính năng đã
có hoặc như kết quả chắc chắn sẽ cải thiện.

\subsection{Lộ trình phát triển có thể kiểm chứng}
Ưu tiên đầu tiên là làm rõ hợp đồng dữ liệu: định nghĩa đơn bị trễ,
cách xử lý đơn hủy, múi giờ và thời điểm có từng thuộc tính. Nếu thay
nhãn hoặc quần thể, cần lưu phiên bản dữ liệu và protocol mới, không
ghép trực tiếp metric cũ với metric mới. Việc xác nhận này quyết định
mô hình đang giải đúng bài toán nghiệp vụ hay chỉ khớp nhãn trong CSV.

Ưu tiên thứ hai là đánh giá tổng quát hóa. Có thể mở rộng mẫu theo ngân
sách, lặp group split với nhiều seed và bổ sung temporal holdout được
giữ độc lập. Kết quả nên có khoảng tin cậy bootstrap theo đơn; các dòng
trong cùng đơn phải đi cùng nhau ở mỗi lần resample. Nếu mục tiêu phục
vụ khách mới, cần thêm thiết kế tách khách hàng thay vì chỉ tách đơn.

Ưu tiên thứ ba là thiết kế tiêu chí vận hành. Cần xác định mức Recall
tối thiểu, số cảnh báo có thể xử lý và chi phí FP/FN. Threshold hoặc
calibration chỉ được học trên validation, không dùng test hiện tại.
Recall Standard Class cần được theo dõi riêng để tránh mô hình cải
thiện điểm trung bình nhưng tiếp tục bỏ sót một nhóm lớn.

Sau cùng mới cân nhắc mở rộng biểu diễn hoặc mô hình: mã hóa chu kỳ
cho biến lịch, đặc trưng vận hành có thời điểm rõ ràng, hay các cấu hình
kiểm soát lớp và độ phức tạp. Mọi phương án cần đối chứng cùng protocol
và có ngân sách tính toán ghi rõ. Không mặc định thêm mô hình phức tạp
sẽ vượt baseline hoặc giải được điểm yếu Standard Class.

\subsection{Điều kiện ứng dụng thực tế}
Nếu các bước trên đạt yêu cầu, hệ thống có thể được thử ở chế độ chỉ
ghi nhận, theo dõi schema, tỷ lệ danh mục lạ, phân phối xác suất và metric
khi nhãn hoàn thiện. Việc đánh giá cần tính độ trễ xuất hiện nhãn, vì
chưa giao xong không đồng nghĩa không trễ. Một cơ chế ghi nhật ký giúp
phân biệt thay đổi mô hình với thay đổi cách vận hành.

Để chứng minh lợi ích thực tế, cần đo thêm tác động của quy trình xử lý
cảnh báo, không chỉ F1. Tỷ lệ cảnh báo đúng, thời gian nhân viên xử lý
và kết quả sau can thiệp phải được thu thập theo kế hoạch phù hợp. Trong
phạm vi hiện tại, đóng góp của đề tài là một quy trình có thể kiểm tra,
các kết quả thực nghiệm được truy vết và nhận diện rõ giới hạn cần
khắc phục trước khi mở rộng.
